%!TEX TS-program = lualatex
%!TEX encoding = UTF-8 Unicode

\documentclass[11pt, addpoints]{exam}

%\printanswers

\usepackage[left=1in,right=1.25in,top=1in]{geometry}     

\usepackage{fontspec}
\def\mainfont{Linux Libertine O}
\setmainfont[Ligatures={TeX, Common}, BoldFont={* Bold}, ItalicFont={* Italic}, Fractions ={On}, Numbers={Proportional}]{\mainfont}
\setsansfont[Scale=MatchLowercase]{Linux Biolinum O} 
\usepackage{microtype}

\usepackage{graphicx}
	\graphicspath{%
	{/Users/goby/Pictures/teach/348/exams/}}%

% Used to get the last two digits of the year for the header below.
\def\short#1{\csname @gobbletwo\expandafter\endcsname\number#1}

\pagestyle{headandfoot}
\firstpageheader{Marine Biology, Exam 4, Sp\short{\year}.}{}{%
	\ifprintanswers\textbf{KEY}\else Name: \enspace \makebox[2.5in]{\hrulefill}\fi}
\runningheader{}{}{\small{pg. \thepage}}
%\runningheadrule

\footer{\iflastpage{\small \rule{0.6in}{0.4pt} / \numpoints\ points}{}}{}{\iflastpage{\small Have a great summer!}{}}%{\makebox[0.75in]{\hrulefill}}
\extrafootheight{-0.5in}

\pointsinmargin
\marginpointname{ pts}
\marginbonuspointname{ pts \textsc{e.c.}}


\usepackage{multicol}
\usepackage[version=4]{mhchem}

%% Remove comment %% to print answer key.
\unframedsolutions
\renewcommand{\solutiontitle}{}
\SolutionEmphasis{\bfseries}

%% Create a True False question format
\newcommand*{\TrueFalse}[1]{%
\ifprintanswers
	\ifthenelse{\equal{#1}{T}}{%
		\textbf{TRUE}\hspace*{14pt}False
	}{
		True\hspace*{14pt}\textbf{FALSE}
	}
\else
	{True}\hspace*{20pt}False
\fi
}
\newlength\TFlengthA
\newlength\TFlengthB
\settowidth\TFlengthA{\hspace*{1.16in}}
\newcommand\TFQuestion[2]{%
	\setlength\TFlengthB{0.98\linewidth}
	\addtolength\TFlengthB{-\TFlengthA}
	\parbox[t]{\TFlengthA}{\TrueFalse{#1}}\parbox[t]{\TFlengthB}{#2}}


%% Create a Matching question format
\newcommand*\Matching[1]{
\ifprintanswers
	\textbf{#1}
\else
	\rule{2in}{0.4pt}
\fi
}
\newlength\matchlena
\newlength\matchlenb
\settowidth\matchlena{\rule{2.1in}{0pt}}
\newcommand\MatchQuestion[2]{%
	\addtocounter{numpoints}{2}
	\setlength\matchlenb{0.98\linewidth}
	\addtolength\matchlenb{-\matchlena}
	\parbox[t]{\matchlena}{\Matching{#1}}\enspace\parbox[t]{\matchlenb}{#2}}


% change indentation for multiple choice
\renewcommand{\choiceshook}{%
	\setlength{\leftmargin}{0.25in}%
	\setlength{\parsep}{3pt}%
}

\newcommand*\AnswerBox[2]{%
    \parbox[t][#1]{0.92\textwidth}{%
    \begin{solution}#2\end{solution}}
    \vspace*{\stretch{1}}
}

\newenvironment{AnswerPage}[1]
    {\begin{minipage}[t][#1]{0.92\textwidth}%
    \begin{solution}}
    {\end{solution}\end{minipage}
    \vspace*{\stretch{1}}}

\newlength{\basespace}
\setlength{\basespace}{5\baselineskip}


\begin{document}

\begin{questions}

\fullwidth{Read each question carefully and completely. Be sure your answer addresses the question being asked. Use complete sentences. Use your words.}

\bonusquestion[1] Favorite type of donut (no promises): \rule{2.5in}{0.4pt}

\question[5]
Mudskippers can absorb O\textsubscript{2} through their skin. Briefly describe another mechanism they use to breathe on exposed mudflats. How do mudskippers use this mechanism to deliver oxygen to their developing eggs?

\AnswerBox{0.2\textheight}{Trap a bubble of air over their gills. Males will deliver this bubble into the burrow and release it over their eggs.}

\question[5]
Compare and contrast the plasticity of the Atlantic Puffin versus the Little Auk. Which species shows more ability to adapt to its changing environment? 

\AnswerBox{0.2\textheight}{Puffins abandoning nests and experiencing low reproductive success when food resources change, whereas little auks adapt to the change in food resources and adjust diving behavior according to sea-ice coverage, benefit from glacial melt and rising temps.}

%\newpage

\question[5]
Briefly describe how predators affect the distribution of prey species (like snails) in a salt marsh during low and high tides.

\AnswerBox{0.2\textheight}{
	Prey species tend to move about in the water to feed during low tide because predators are restricted in their movement. Prey move up on the stems above the water line during high tide to avoid the predators.}

\newpage


\question[5] \label{q:carbonate_buffering_system}
Write the equilibrium equation for the carbonate buffering system in
marine ecosystems, with correct valence for ions. Use this equation to explain how the increase in
atmospheric \ce{CO2} increases ocean acidity. %Circle the hydrogen ion 
%that contributes to most increase in acidity.

%\textbf{30 points i based on carbonate buffering system. This seems like
%	too many points. Replace one of the questions with something about
%	biological response to acidification.}

\begin{AnswerPage}{0.2\textheight}%

\ce{CO2 + H2O <=> H2CO3 <=> H+ + HCO3- <=> 2H+ + CO3^2-} (2 points)

\vspace{\baselineskip}

As \ce{CO2} is added to seawater, the equation drives to the right to achieve equilibrium. This increases the amount of \ce{H+} dissolved in water. pH is a measure of free \ce{H+}.  More free \ce{H+} means lower pH, meaning higher acidity. (3 points)

%\vspace{\baselineskip}

%2 points for the entire equilibrium equation, 3 points for the explanation.% Grade as 1 point per molecule. Deduct 1 point for each incorrect superscript or subscript. Deduct 1 point if not shown as equilibrium equation. 1 point for the circle. 4 points for the explanation.

\end{AnswerPage}


\question[5] \label{q:calcium_carbonate_equilibrium}
Write the equilibrium equation for \ce{CaCO3} chemistry, again with correct valence.

\AnswerBox{0.1\textheight}{%
\ce{CaCO3 <=> Ca^2+ + CO3^2-}   \\
Grade as 1 point per molecule. Deduct 1 point for each incorrect superscript or subscript. Deduct 1 point if not shown as equilibrium equation.
}

\question[10]
Explain how the increased acidity (from the equilibrium equation in
question~\ref{q:carbonate_buffering_system}) interacts with \ce{CO3^2-} (from the equilibrium equation in question~\ref{q:calcium_carbonate_equilibrium})
to reduce \ce{CaCO3} availability to marine animals. Include
the proper chemical equation as part of your answer 

%\textbf{Write 15 pt carbonate buffering system equation. Explain how increased CO2 causes decreased CO2 availability.}

\begin{AnswerPage}{0.4\textheight}%

\ce{H+ + CO3^2- -> HCO3-} \qquad (2 pts)

\vspace{\baselineskip}

The hydrogen ion from the carbonate buffering system combines with the carbonate ion from the calcium carbonate system to form bicarbonate. (3 pts) 

This prevents the carbonate ion from being available to organisms that produce calcium carbonate skeletons. (3 pts)

\end{AnswerPage}

\question[5]
What is the minimum value of \Omega\ necessary for carbonate skeleton
formation? What is the optimum value?

\AnswerBox{0.1\textheight}{%
Minimum $>$ 1 for shell formation. (2 pts)\\
Optimum $\ge$ \textasciitilde4.0 (2 pts)
}

\question[5]
Describe how ocean acidification might affect the ability of juvenile coral reef fishes (e.g., clownfish) to find the proper habitat to recruit into.

\AnswerBox{0.2\textheight}{%
	Tended to go to wrong habitat under acidic conditions. Also tended to be unable to identify predatory fishes, decreasing chance of surviving.
}


\newpage

%\question[5]
% MSY is based on a population logistic growth curve. Draw the logistic
%growth curve, and then draw an arrow to indicate the point along the
%growth curve that provides the theoretical maximum sustainable yield.
%
%\AnswerBox{1\basespace}{%
%2 pts for the curve. \\
%2 pts: Arrow pointing to the steepest part of the curve.
%}


\newpage

\fullwidth{\centering\includegraphics[width=6.5in]{Exam4_Fig1a}}

\question[15]
Use the figure above to clearly and thoroughly explain how increased ocean acidity
will alter connectivity and fecundity across a large coral reef. You must include coral
recruitment and growth (indicated by the black arrow), distance between patches, and coral fecundity (reproductive output of corals). Be sure to include both positive (dashed lines) and negative (solid lines) relationships among these factors. Do not include unnecessary factors and relationships. 

\begin{solution}
Grade as 2 points per item. 

Acidification will negatively affect coral recruitment and growth. Fewer new coral recruits will be added to the population and growth of existing corals will be inhibited.  

Reduced coral growth has a negative relationship with bleached and dead corals. As coral growth is reduced, coral bleaching and death increases.  

Coral bleaching and death has a positive relationship with distance between factors. Thus, as coral bleaching/death increases, the distance between the patches also increases.

Patch distance has a negative influence on coral recruitment and growth through a reduction in connectivity.

Coral recruitment and growth has a positive influence on diversity of fecund corals. as coral recruitment decreases so will the diversity of fecund corals. 

In addition, bleached and dead corals has a negative relationship with diversity of fecund corals. As coral loss increases, the diversity of fecund corals decreases. 

Diversity of fecund corals has a positive relationship with reef scale coral fecundity. As diversity of fecund corals decrease in one patch, so does fecundity across the entire coral reef (reef scale).

Reef scale fecundity has a positive influence on coral recruitment and growth. As reef scale fecundity decreases, the larval supply will decrease, decreasing coral recruitment.

\end{solution}

\newpage

\question[5]
List in order from ocean to land the four zones found in a typical New England salt marsh. You may use common or scientific names but scientific names make a greater impression. 

\AnswerBox{0.1\textheight}{Cordgrass zone, salt hay zone, black rush zone, marsh elder zone.}

\question[10]
Describe the \emph{specific} biotic and abiotic interactions between Cordgrass (\textit{Spartina alterniflora}) and Salt Hay (\textit{S. patens}) that determine their distribution in a salt marsh. Include in your answer which species is found in the low marsh and which is found in the high marsh. 

\begin{AnswerPage}{0.6\textheight}%
	1. Patens (high marsh) is better competitor for space but it cannot tolerate wet, low O\textsubscript{2} of sediment in low marsh.
	
	2. Alterniflora cannot compete with patens for space but can tolerate wet, low O\textsubscript{2} of sediment in low marsh.
\end{AnswerPage}


\end{questions}

\end{document}